% TOP_PROBLEM: {"id": 1349, "title": "GCH and Diamond Principle", "category_id": 10, "category": {"id": 10, "name": "set_theory", "display_name": "Set Theory", "description": "Foundations of mathematics, infinite sets, and cardinality.", "slug": "set-theory", "order_index": 10, "created_at": "2026-07-31T15:26:25.671Z"}, "status": "open"}
% FIRST_POSTED: 2026-10-01
% ATTEMPT_STATUS: unresolved
% Imported from: unsolvedmath_additional_500.tex; UnsolvedMath ID 1349
% !TeX program = lualatex
% Compile twice: lualatex 1349.tex
% Self-contained: no external bibliography, images, or \input files.
% Numeric section labels and visible numbers are original UnsolvedMath IDs.
\documentclass[11pt,a4paper]{article}
\usepackage[margin=24mm,headheight=14pt]{geometry}
\usepackage{fontspec}
\setmainfont{Latin Modern Roman}
\setsansfont{Latin Modern Sans}
\setmonofont{Latin Modern Mono}
\usepackage{amsmath,amssymb,mathtools}
\usepackage{unicode-math}
\setmathfont{Latin Modern Math}
\usepackage{microtype}
\usepackage{enumitem,xurl,needspace}
\usepackage[dvipsnames]{xcolor}
\usepackage{hyperref}
\hypersetup{unicode=true,colorlinks=true,linkcolor=MidnightBlue,urlcolor=MidnightBlue,
 pdftitle={UnsolvedMath Problem 1349},pdfauthor={Source-annotated compilation},bookmarksnumbered=true}
\usepackage{bookmark,tocloft,titlesec,fancyhdr}
\setlength{\cftsecnumwidth}{4.2em}
\cftsetpnumwidth{2.5em}
\cftsetrmarg{4em}
\titleformat{\section}{\Large\bfseries\raggedright}{\thesection}{0.7em}{}
\pagestyle{fancy}\fancyhf{}
\fancyhead[L]{\small\sffamily Additional UnsolvedMath statements}
\fancyhead[R]{\small Original catalogue IDs}
\fancyfoot[C]{\thepage}
\setlength{\parindent}{0pt}
\setlength{\parskip}{5pt plus 1pt}
\setlength{\emergencystretch}{3em}
\setcounter{tocdepth}{1}\setcounter{secnumdepth}{1}
\setlist[itemize]{leftmargin=3em,itemsep=4pt,topsep=4pt,parsep=0pt}
\allowdisplaybreaks
\newcommand{\N}{\mathbb{N}}
\newcommand{\Z}{\mathbb{Z}}
\newcommand{\Q}{\mathbb{Q}}
\newcommand{\R}{\mathbb{R}}
\newcommand{\C}{\mathbb{C}}
\newcommand{\F}{\mathbb{F}}
\newcommand{\A}{\mathbb{A}}
\newcommand{\PP}{\mathbb{P}}
\newcommand{\E}{\mathbb{E}}
\newcommand{\eps}{\varepsilon}
\newcommand{\dd}{\,\mathrm d}
\newcommand{\sref}[2]{\hyperlink{source:#1:#2}{[S#2]}}
\newif\ifshowcatalogue
\showcataloguetrue
\newcommand{\cataloguescope}[1]{\ifshowcatalogue
 \paragraph{Statement recorded in the supplied source.}
 #1\fi}
\begin{document}
% UnsolvedMath ID 1349; source catalogue code ST-003

\setcounter{section}{1348}

\section{Diamond at the Cofinality of a Singular Cardinal}\label{1349}

{\small\sffamily UnsolvedMath ID 1349 · Set Theory}

\subsection{Definitions and mathematical statement}

\paragraph{Shared notation.}Unless a section says otherwise, $\N=\{1,2,\ldots\}$,
$\Z$ is the ring of integers, $\Q,\R,\C$ are the rational, real, and complex
fields, $[N]=\{1,\ldots,\lfloor N\rfloor\}$, and $\log$ is the natural
logarithm. For real functions of a parameter tending to infinity,
$f\ll g$ and $f=O(g)$ mean $|f|\leq Cg$ eventually for some fixed $C$;
$f\gg g$ reverses this comparison; $f=o(g)$ means $f/g\to0$;
$f\sim g$ means $f/g\to1$; and $f\asymp g$ means both order bounds.
A subscript on $O$ or $\ll$ specifies allowed constant dependence.
The expression $n^{o(1)}$ in an upper bound means growth smaller than
$n^\epsilon$ for each fixed $\epsilon>0$. Local definitions govern any
exception, finite-field convention, or use of the same symbol for another
object. An asymptotic claim is not automatically an effective algorithm.



For singular $\lambda$, put $S=E_{\operatorname{cf}(\lambda)}^{\lambda^+}=\{\alpha<\lambda^+:\operatorname{cf}(\alpha)=\operatorname{cf}(\lambda)\}$. The principle $\diamondsuit(S)$ is a sequence $(A_\alpha)_{\alpha\in S}$, $A_\alpha\subseteq\alpha$, such that for every $A\subseteq\lambda^+$, the set $\{\alpha\in S:A\cap\alpha=A_\alpha\}$ is stationary. Stationary means meeting every closed unbounded subset of $\lambda^+$. \sref{1349}{1} \sref{1349}{2}

\paragraph{Statement recorded in the supplied source.}
Does the generalized continuum hypothesis entail the diamond principle $\diamondsuit(E_{\text{cf}(\lambda)}^{\lambda^+})$ for every singular cardinal $\lambda$? \sref{1349}{1}

\subsection{Short English statement}

Does the generalized continuum hypothesis supply diamond prediction on the cofinality class associated with every singular cardinal?

\subsection{Research attempt: the necessary arithmetic and stationary guessing obstruction}
\textbf{Status.} We prove consequences that any proposed diamond sequence must satisfy and explain why the cardinal count supplied by GCH does not construct that sequence. The implication in the question remains unresolved. The primary survey [R1] treats diamond and weaker guessing principles at successors and records the distinction used below; none of its deeper equivalences is assumed in the proofs.

\subsubsection{1. The prescribed cofinality set really is stationary}
Write $\theta=\operatorname{cf}(\lambda)$ and $\kappa=\lambda^+$. For any club $C\subseteq\kappa$, recursively choose a strictly increasing continuous sequence $(\alpha_i)_{i<\theta}$ in $C$. Its supremum $\alpha$ is below $\kappa$ by regularity of $\kappa$, and belongs to $C$ by closure. It has cofinality exactly $\theta$: it has a cofinal sequence of that length, while a shorter cofinal sequence would give a cofinal subset of the regular ordinal $\theta$ of smaller size by looking at first indices dominating its entries. Therefore $\alpha\in C\cap E_\theta^\kappa$. This verifies the ambient stationary set on which the target operates.

\subsubsection{2. Diamond forces GCH at the predecessor}
Assume $(A_\alpha)_{\alpha\in S}$ witnesses $\diamondsuit(S)$, where $S=E_\theta^\kappa$. For each $B\subseteq\lambda$, regard $B$ as a subset of $\kappa$. There is an $\alpha\in S$ with $\alpha>\lambda$ and $A_\alpha=B$: the stationary guessing set meets the club tail above $\lambda$, and then $B\cap\alpha=B$. Hence the map
\[\alpha\longmapsto A_\alpha\cap\lambda\]
from $S$ onto a family containing every subset of $\lambda$ shows $2^\lambda\leq\kappa$. Cantor's theorem gives the reverse lower bound $2^\lambda\geq\lambda^+$, so
\[\diamondsuit(E_\theta^{\lambda^+})\ \Longrightarrow\ 2^\lambda=\lambda^+.
\tag{1}\]
This proves necessity of the local arithmetic condition. It is not the converse asked for, and the arrow in (1) cannot be reversed by counting alone.

\subsubsection{3. Stationary success is a simultaneous requirement}
The actual diamond quantifiers are
\[\text{for every }A\subseteq\kappa\text{ and every club }C\subseteq\kappa,
\text{ there is }\alpha\in C\cap S\text{ with }A_\alpha=A\cap\alpha.\tag{2}\]
Under GCH at $\lambda$ one may enumerate $\mathcal P(\lambda)$ in $\kappa$ steps. This addresses the number of small subsets, but not the $2^\kappa$ possible full targets in (2), and it supplies no coherence among their restrictions to different ordinals. Nor does unbounded guessing imply stationary guessing: the successor ordinals are unbounded in $\kappa$ and disjoint from the club of limit ordinals.

A further direct consequence illustrates the force of (2). For each $\xi<\kappa$ consider
\[S_\xi=\{\alpha\in S:\alpha>\xi, A_\alpha=\{\xi\}\}.
\]
Applying diamond to the singleton $\{\xi\}$ shows $S_\xi$ is stationary. These sets are pairwise disjoint, since a single $A_\alpha$ cannot equal two different singletons. Thus one sequence allocates stationary, rather than merely nonempty, success sets even to these $\kappa$ elementary targets. This consequence alone is not a characterization of diamond.

\subsubsection{4. Adversarial checks and remaining gap}
The proof of (1) restricts to $\alpha>\lambda$; without that restriction equality to the whole subset $B$ would not follow. In the stationarity proof $\theta$ is regular because it is a cofinality, and $\kappa$ is regular because it is a successor cardinal in ZFC. These are the specific regularity hypotheses, not an assumption that the singular cardinal $\lambda$ is regular. The proposed enumeration argument fails precisely at the simultaneous club quantifier in (2). No sequence satisfying that quantifier is constructed, and no model refuting its existence under GCH is supplied. That stationary coherence problem is the remaining content of the catalogue question.

\subsection{Sources}

{\small

\begin{itemize}

\item[\hypertarget{source:1349:1}{S1}] ULAM.AI, \emph{UnsolvedMath}, supplied \texttt{problems.json}, record 1349 (ST-003), statement and background. The embedded literature-review date is 2026-08-17; that is the archive’s review date, not a fresh certification. Dataset landing page: \url{https://huggingface.co/datasets/ulamai/UnsolvedMath}.

\item[\hypertarget{source:1349:2}{S2}] T. Jech, Set Theory, 3rd millennium edition (2003). \url{https://doi.org/10.1007/3-540-44761-X}. Bibliographic information imported from the supplied record.

\item[R1] A. Rinot, Jensen’s diamond principle and its relatives, Contemporary Mathematics 533 (2011); inspected primary survey record.
\url{https://arxiv.org/abs/0911.2151}.
\end{itemize}

}

\label{end:1349}
\end{document}
